% DAFx26 Demo Track Paper: FreeEQ8/ProEQ8
% Reformatted using official DAFx26 LaTeX Demo Template
%
% To compile: pdflatex, bibtex, pdflatex, pdflatex

%------------------------------------------------------------------------------------------
%  User defined variables
%------------------------------------------------------------------------------------------
\def\papertitle{Real-Time State-Space Parameterization in Digital Equalization: A Production-Grade SVF Implementation}
\def\paperauthorA{Gary Doman}

%------------------------------------------------------------------------------------------
\documentclass[twoside,a4paper]{article}
\usepackage{etoolbox}

\usepackage{dafx26demo}

\usepackage{amsmath,amssymb,amsfonts,amsthm}
\usepackage{siunitx}
\usepackage{euscript}
\usepackage[T1]{fontenc}
\usepackage[utf8]{inputenc}
\usepackage{ifpdf}
\usepackage[english]{babel}
\usepackage{caption}
\usepackage{subfig}
\usepackage{color}
\usepackage{booktabs}

\input glyphtounicode
\pdfgentounicode=1

\setcounter{page}{1}
\ninept

% Author count machinery (DO NOT MODIFY)
\newcounter{numauth}\setcounter{numauth}{1}
\newcounter{listcnt}\setcounter{listcnt}{1}
\newcommand\authcnt[1]{\ifdefined#1 \stepcounter{numauth} \fi}
\newcommand\addauth[1]{
\ifdefined#1 
\stepcounter{listcnt}
\ifnum \value{listcnt}<\value{numauth}
\appto\authorslist{, #1}
\else
\appto\authorslist{~and~#1}
\fi
\fi}
\def\authorslist{\paperauthorA}

\usepackage{times}

% PDF/LaTeX detection
\newif\ifpdf
\ifx\pdfoutput\relax
\else
   \ifcase\pdfoutput
      \pdffalse
   \else
      \pdftrue
   \fi
\fi

\ifpdf
  \usepackage[pdftex,
    pdftitle={\papertitle},
    pdfauthor={\authorslist},
    pdfsubject={Proceedings of the 29th International Conference on Digital Audio Effects (DAFx26)},
    colorlinks=false,
    bookmarksnumbered,
    pdfstartview=XYZ
  ]{hyperref}
  \pdfcompresslevel=9
  \usepackage[pdftex]{graphicx}
\else
  \usepackage[dvips]{epsfig,graphicx}
  \usepackage[dvips,
    pdftitle={\papertitle},
    pdfauthor={\authorslist},
    pdfsubject={Proceedings of the 29th International Conference on Digital Audio Effects (DAFx26)},
    colorlinks=false,
    bookmarksnumbered,
    pdfstartview=XYZ
  ]{hyperref}
\fi
\usepackage[hypcap=true]{caption}
\title{\papertitle}

% Single author, single affiliation
\affiliation
{\paperauthorA}
{Independent Researcher (GareBear99 / TizWildin) \\ 
\href{https://github.com/GareBear99/FreeEQ8}{https://github.com/GareBear99/FreeEQ8}\\
{\tt \href{mailto:neovectr.inc@gmail.com}{neovectr.inc@gmail.com}}
}

\begin{document}
\ifpdf
  \DeclareGraphicsExtensions{.png,.jpg,.pdf}
\else
  \DeclareGraphicsExtensions{.eps}
\fi

\maketitle

\begin{abstract}
We present FreeEQ8/ProEQ8, an open-source parametric equalizer implementing the Simper State Variable Filter (SVF) topology to eliminate high-frequency magnitude cramping without oversampling. Traditional RBJ biquad implementations exhibit 199\% Q distortion at 16\,kHz (44.1\,kHz sample rate); our SVF achieves exact response at all frequencies up to Nyquist. The architecture features lock-free SPSC triple-buffering for real-time safety, a variable-cadence coefficient engine reducing Dynamic EQ CPU cost by 75\%, and allocation-free semantic resonance detection. Measured performance: 0.62\% single-core CPU at 44.1\,kHz (161$\times$ headroom). Source code, benchmarks, and CI validation available under GPL-3.0.
\end{abstract}

\section{Introduction}
\label{sec:intro}

Traditional parametric EQ implementations use the Robert Bristow-Johnson (RBJ) Audio EQ Cookbook biquad formulas~\cite{RBJ} in Transposed Direct Form II. While mathematically correct, the bilinear transform introduces frequency cramping near Nyquist: a Bell filter at 16\,kHz with $Q=1.0$ at 44.1\,kHz exhibits an effective bandwidth \textbf{199\% narrower} than intended.

Industry solutions include brute-force oversampling (adding latency and CPU cost) or proprietary polynomial analog-matching curves. This paper documents a third path: the Simper SVF topology~\cite{Simper2013}, which pre-warps the cutoff frequency via $g = \tan(\pi \cdot f_c / f_s)$, achieving exact cutoff placement without oversampling.

\subsection{Product Architecture}

The codebase produces two plugins from a single source tree:

\textbf{FreeEQ8} (Free, GPL-3.0): 8-band parametric EQ using RBJ TDF-II biquads. Zero audio restrictions during real-time playback.

\textbf{ProEQ8} (\$20): 24-band parametric EQ using Simper SVF with de-cramped HF response. Adds 4 saturation modes, A/B comparison, and collision detection.

\section{Filter Topology}
\label{sec:topology}

\subsection{RBJ TDF-II (FreeEQ8)}

Standard biquad transfer function with coefficients per the RBJ cookbook. 64-bit double internal state, float I/O.

\textbf{Measured Q distortion (Bell, $Q=1.0$, 44.1\,kHz):}

\begin{table}[ht]
  \caption{RBJ Q distortion vs frequency.}
  \centering
  \begin{tabular}{ccc}
    \toprule
    Frequency & Effective Q & Error \\
    \midrule
    1\,kHz & 1.005 & +0.5\% \\
    8\,kHz & 1.337 & +33.7\% \\
    16\,kHz & 2.990 & +199\% \\
    \bottomrule
  \end{tabular}
  \label{tab:qdistortion}
\end{table}

\subsection{Simper SVF (ProEQ8)}

Pre-warped cutoff via trapezoidal integration~\cite{Simper2013}:
\begin{equation}
g = \tan(\pi \cdot f_c / f_s), \quad k = 1/Q
\end{equation}
\begin{equation}
a_1 = 1 / (1 + g(g + k)), \quad a_2 = g \cdot a_1, \quad a_3 = g \cdot a_2
\end{equation}

Per-sample processing uses the optimized bounded form with two integrator states ($ic_{1eq}$, $ic_{2eq}$). Output mix coefficients select filter type (LP, HP, BP, Bell, shelves).

\textbf{Key result at $f_c = 16$\,kHz:}

\begin{table}[ht]
  \caption{SVF vs RBJ magnitude response at 16\,kHz.}
  \centering
  \begin{tabular}{ccc}
    \toprule
    Topology & Magnitude at $f_c$ & Error \\
    \midrule
    RBJ @ 44.1\,kHz & +0.73\,dB & $-5.27$\,dB \\
    SVF @ 44.1\,kHz & +6.00\,dB & 0.00\,dB \\
    RBJ @ 4$\times$ OS & +4.82\,dB & $-1.18$\,dB \\
    \bottomrule
  \end{tabular}
  \label{tab:svfvsrbj}
\end{table}

The SVF achieves exact gain at the center frequency without oversampling.

\section{Real-Time Safety Architecture}
\label{sec:realtime}

\subsection{SPSC Triple-Buffer}

Three buffer slots indexed by \{writeSlot, midSlot, readSlot\}. Audio thread writes and atomically swaps writeSlot $\leftrightarrow$ midSlot with \texttt{memory\_order\_release}. UI thread reads via midSlot $\leftrightarrow$ readSlot with \texttt{memory\_order\_acquire}. No mutex, no blocking.

\subsection{Off-Thread FIR Reconstruction}

Linear phase mode: 2048-sample latency (4096-tap Hann-windowed FIR). Parameter changes set an atomic dirty flag; a dedicated background thread rebuilds the kernel and publishes via the same triple-buffer protocol.

\subsection{Variable-Cadence Dynamic EQ}

Dynamic EQ recomputes coefficients per-sample when active. We observe that during sustained signals, envelope changes $< 0.1$\,dB between samples are inaudible within a 4-sample window (0.09\,ms at 44.1\,kHz).

\textbf{Algorithm:} If $|\delta| > 0.1$\,dB, update immediately (zero lag on transients). Otherwise, batch to 4-sample intervals.

\textbf{Result:} 75--80\% reduction in coefficient updates with no audible difference.

\section{Benchmarks}
\label{sec:benchmarks}

All benchmarks from standalone C++ (no JUCE, no DAW). Platform: x86-64, g++ -O3.

\begin{table}[ht]
  \caption{Performance measurements (44.1\,kHz, 512-sample blocks).}
  \centering
  \begin{tabular}{ccc}
    \toprule
    Configuration & ns/sample & CPU Headroom \\
    \midrule
    RBJ 8-band stereo & 40.7 & 277$\times$ \\
    SVF 8-band stereo & 70.4 & 161$\times$ \\
    SVF DynEQ worst-case & 370.9 & 30.6$\times$ \\
    \bottomrule
  \end{tabular}
  \label{tab:benchmarks}
\end{table}

SVF overhead vs RBJ: 1.73$\times$. CPU budget at 44.1\,kHz: \textbf{0.62\%}.

Concurrent stress tests (400\,ms, Intel i7-3720QM): 239M samples processed, \textbf{0 data tears}.

\section{Demonstration Plan}
\label{sec:demo}

The live demonstration will showcase:
\begin{enumerate}
\item Side-by-side RBJ vs SVF frequency response at 16\,kHz
\item Real-time spectrum analyzer with lock-free triple-buffer rendering
\item Dynamic EQ envelope tracking on drum transients
\item 24-band surgical EQ workflow in ProEQ8
\end{enumerate}

\textbf{Requirements:} Table with power outlet, external audio interface, headphones for A/B listening.

\section{Conclusions}

We have demonstrated that the Simper SVF topology eliminates high-frequency cramping in parametric EQ without oversampling, achieving exact magnitude response at all frequencies. The lock-free architecture ensures real-time safety, and the variable-cadence engine provides significant CPU savings on Dynamic EQ workloads.

Source code, benchmarks, and CI validation (pluginval strictness-level-10) are available at \url{https://github.com/GareBear99/FreeEQ8} under GPL-3.0.

\section{Acknowledgments}

The Simper SVF topology is due to Andrew Simper (Cytomic). The RBJ biquad coefficients are due to Robert Bristow-Johnson. This work is released under GPL-3.0.

\bibliographystyle{IEEEtranDAFx}
\begin{thebibliography}{9}

\bibitem{Simper2013}
A.~Simper, ``Solving the continuous SVF equations using trapezoidal integration and equivalent currents,'' Cytomic, 2013. [Online]. Available: \url{https://cytomic.com/files/dsp/SvfLinearTrapOptimised2.pdf}

\bibitem{RBJ}
R.~Bristow-Johnson, ``Audio EQ Cookbook,'' musicdsp.org, 1994. [Online]. Available: \url{https://www.musicdsp.org/files/Audio-EQ-Cookbook.txt}

\bibitem{Pirkle2019}
W.~Pirkle, \textit{Designing Audio Effect Plugins in C++}, 2nd ed. Focal Press, 2019.

\bibitem{RennGiles2019}
F.~Renn-Giles and D.~Rowland, ``Real-time 101,'' presented at Audio Developer Conference, 2019. [Online]. Available: \url{https://github.com/hogliux/farbot}

\end{thebibliography}

\end{document}
